\chapter{Introduzione}
\label{cha:Introduzione}

%Data analysis is the process of collecting, cleaning, and transforming data to obtain insights to help make better and informed decisions. In our ever-growing, data-driven world, this is a must for companies of all sizes to solve everyday business problems. Each company has its own team, processes, and tools for data analysis projects.


\section{Contesto storico e turistico della località}\label{par:Contesto_storico}

Cortina d’Ampezzo è una delle località montane più rinomate nel panorama turistico invernale europeo ed è l'area studiata all'interno di questo elaborato. Situata nel cuore delle Dolomiti Bellunesi, ha costruito la propria identità turistica a partire dalla fine del XIX secolo, diventando progressivamente una meta d’eccellenza per il turismo alpino internazionale. La consacrazione definitiva avvenne nel \textbf{1956}, quando Cortina ospitò la prima edizione italiana dei Giochi Olimpici Invernali, guadagnandosi un posto di rilievo tra le più importanti destinazioni sciistiche del mondo \cite{In_focus_HVS}.

Oggi, la località rappresenta un polo turistico di primaria importanza, sia per il mercato italiano, di cui nel 2024 è stata la punta di diamante \cite{Cronaca2024VacanzeNeve}, sia per quello internazionale. Il turismo costituisce da decenni il principale motore economico del territorio ampezzano, sostenuto da infrastrutture moderne, eventi sportivi di rilievo mondiale e grandi produzioni cinematografiche (\cite{Only_for_your_eyes} \cite{Vacanze_di_Natale} \cite{Vacanze_di_Natale_a_Cortina}). Gli eventi sportivi ospitati, come la Coppa del Mondo e i Mondiali del 2021 di sci alpino o la Coppa del Mondo di Snowboard Cross, si svolgono annualmente negli impianti sciistici della Tofana, rafforzando l'immagine di Cortina come meta turistica internazionale.\\

Il comprensorio sciistico cortinese, e in particolare la Tofana, rappresenta il principale sfondo geografico di riferimento per questo elaborato. La dimensione principale dei dati analizzati è costituita dai passaggi ai tornelli degli impianti sciistici da parte dei turisti, forniti dalla società \textbf{Freccia nel Cielo}, che gestisce una parte significativa, sebbene non la totalità, degli impianti di risalita situati in quest'area. L'azienda, la cui storia sarà approfondita nel capitolo \ref{cap:Storia_FNC}, ha condiviso i dati per finalità di ricerca, rendendo possibile un'analisi quantitativa e qualitativa delle dinamiche di affluenza e dell'utilizzo degli impianti.\\

Questi dati consentono di stimare il volume di sciatori presenti sulle piste e, integrati con la conoscenza del contesto territoriale e storico, offrono chiavi di lettura più accurate delle dinamiche di affluenza turistica.

È opportuno precisare che la Tofana non costituisce l’unica area sciistica della località ampezzana: il comprensorio comprende anche \textbf{Faloria–Cristallo}, \textbf{Cinque Torri–Col Gallina–Lagazuoi}, \textbf{San Vito di Cadore} e \textbf{Misurina-Auronzo di Cadore}, tutte accessibili con uno skipass unico di valle. Tuttavia, la Tofana registra il flusso di sciatori più elevato grazie alla vicinanza al centro abitato, alla maggiore estensione e ai collegamenti con le altre aree, oltre a una distribuzione omogenea dei livelli di difficoltà delle piste. Questi fattori ne fanno l’area più rappresentativa per l’analisi svolta.


\section{Motivazioni dello studio}
Lo sviluppo turistico della valle ampezzana ha determinato una crescente pressione sulle infrastrutture e sui servizi, richiedendo una pianificazione efficiente per mantenere competitività e qualità. Gli impianti sciistici, gli albergatori e gli operatori locali devono gestire un’elevata affluenza nei periodi di alta stagione, in particolare durante le festività natalizie e di Capodanno, quando si registra il picco annuale del tasso di occupazione. Nel 2024, questo picco ha raggiunto circa il 90\% della capacità disponibile \cite{Andreazza2024CortinaTourismo}, evidenziando la forte pressione sul sistema turistico.\\

Considerando l’elevata occupazione nei periodi festivi e le dimensioni effettive del comune ampezzano, che conta meno di 5.500 abitanti (dato 2024 \cite{cortina_wiki}),risulta fondamentale una gestione organizzativa precisa e scalabile delle piste da sci per garantire la qualità dell’esperienza turistica e ottimizzare le risorse.


Questo quadro di sviluppo turistico e di domanda di infrastrutture costituisce la cornice empirica su cui si fonda l’analisi presentata in questo lavoro: l’obiettivo è analizzare i dati relativi ai passaggi sugli impianti di Cortina d'Ampezzo, nelle strutture possedute dalla precedemente citata Freccia nel cielo,  e compararli nei periodi di affluenza , con ulteriori dimensioni dei dati quali, informazioni meteorologiche. 
Lo scopo ultimo è ottenere dei risultati che possano essere significativi o inaspettati, identificare pattern comuni e  fornire dati utili alla località negli ambiti di gestione delle risorse, organizzativo e turistico. \\

\section{Prospettive olimpiche}
La “Perla delle Dolomiti” è stata selezionata insieme a Milano come sede delle Olimpiadi Invernali 2026, a settant’anni dalla prima edizione \cite{Milano_Cortina_olympics}. L’annuncio ufficiale, avvenuto il 24 giugno 2019 \cite{Olympic_nominee}, è stato il risultato di un processo di selezione basato su criteri di sostenibilità, accessibilità e capacità infrastrutturale.\\
Questa designazione ha già generato un importante effetto mediatico e promozionale, contribuendo a modellare le percezioni di turisti e stakeholder sul potenziale della località. Sebbene l’evento olimpico non rientri direttamente nel periodo di osservazione dello studio, la sua imminenza può aver influenzato indirettamente i flussi turistici attraverso un aumento della visibilità e delle prenotazioni.
Un’estensione futura dell’analisi potrà includere il confronto tra periodi pre e post-olimpici per valutare l’impatto a medio e lungo termine su affluenza, utilizzo degli impianti e dinamiche di prezzo.


\section{Metodologia di ricerca}
L'integrazione tra dati operativi e dati metereologici rappresenta il nucleo metodologico di questo studio, poiché consente di estrarre informazioni preziose per la gestione quotidiana e a lungo termine, da un punto di vista strategico,  degli impianti sciistici.

Le informazioni ricavate verranno presentate  sotto forma di visualizzazioni interattive e permetteranno un supporto al processo decisionale più consapevole: identificare i picchi di domanda, prevedere le giornate critiche, ottimizzare la distribuzione dei carichi di lavoro e gestire le risorse in maniera più efficace, sono alcuni degli esempi applicativi reali di questo studio.

Per conseguire tali obiettivi, la metodologia adottata integra fonti di dati quantitativi e strumenti informatici per l’analisi multidimensionale. I dataset considerati comprendono i registri dei passaggi sugli impianti sciistici, associati a timestamp giornalieri, e i dati meteorologici riferiti alle stesse stagioni invernali oggetto di studio. A tali informazioni sono state affiancate variabili di contesto relative a festività ed eventi locali, con l’obiettivo di arricchire l’analisi e ampliare la capacità esplicativa del modello rispetto ai fenomeni di affluenza. 

L’intervallo temporale dei dato copre cinque stagioni (2019/20–2023/24), includendo sia il periodo di sospensione delle attività dovuto alla pandemia di COVID-19 sia la successiva ripresa, per individuare eventuali discontinuità e trend strutturali.\\

La pipeline metodologica sviluppata comprende:

\begin{enumerate}
    \item Pulizia e normalizzazione dei dati grezzi provenienti da aziende e sistemi meteo.
    \item Segmentazione dei dati per stagioni invernali, selezionando le variabili più significative allo studio.
    \item Modellazione della base dati in uno schema multidimensionale per supportare interrogazioni complesse e viste aggregabili.
    \item Rappresentazione dei risultati tramite dashboard interattive e grafici dinamici.
\end{enumerate}

 Su tali visualizzazioni vengono applicate operazioni tipiche dell’analisi OLAP (Online Analytical Processing), quali:

\begin{itemize}
    \item \textbf{Roll-Up}: operazione che consente di aumentare il livello di astrazione dei dati aggregandoli su dimensioni gerarchiche (ad esempio, da giorno a settimana o mese), permettendo analisi sintetiche e di alto livello vedendo i dati più aggregati.
    \item \textbf{Drill-down}: operazione complemetare al roll-up, permette di esplorare i dati più nel dettaglio (da mese a singolo giorno per esempio), favorendo analisi più mirate nello specifico
    \item \textbf{Slicing e dicing}: tecniche che permettono di isolare e visualizzare specifici sottoinsiemi di dati. Slicing seleziona una singola porzione di dati lungo una dimensione, mentre dicing consente di creare viste multidimensionali più complesse su più dimensioni contemporaneamente.
\end{itemize}

Queste procedure permettono di trasformare i dati grezzi in conoscenza operativa, mettendo in evidenza pattern, variazioni festive e stagionali e fattori esterni che possono alterare l’affluenza. Tali informazioni possono essere sfruttate dagli operatori per anticipare picchi di presenze, ottimizzare la pianificazione delle risorse e modulare l'offerta in relazione ad eventi o condizioni meteo, contribuendo a una gestione più efficiente e strategica della località turistica.\\

L'approccio descritto, applicato agli impianti sciistici della Tofana, può essere replicato anche in altri comprensori montani, offrendo uno strumento flessibile e adattabile per la gestione delle infrastrutture turistiche basato sui dati.


\section{Struttura della tesi}

La struttura dell'elaborato è articolato nei seguenti capitoli

\begin{enumerate}
    \item \textbf{Introduzione}, che presenta il contesto dello studio, gli obiettivi della ricerca e la metodologia adottata.
    \item \textbf{Fonti e caratteristiche dei dati}, dove vengono descritti gli enti che hanno fornito i dati e un po della loro storia, insieme a una presentazione delle principali caratteristiche dei dataset popolati con dati esemplificativi.
    \item \textbf{Operazioni sui dati e modellazione relazionale}, che illustra il processo operativo di preparazione dei dati, dalle operazioni di pulizia e pre–elaborazione alla definizione e documentazione del modello relazionale utilizzato come base per l’analisi.
    \item \textbf{Analisi e visualizzazione dei dati}, in cui vengono presentate le elaborazioni condotte tramite strumenti di data visualization. Le analisi sono guidate da domande di ricerca specifiche, utili a evidenziare i fenomeni osservati e a interpretare i risultati ottenuti.
    \item \textbf{Conclusioni}, dove vengono discussi i risultati principali, le implicazioni dello studio e possibili sviluppi futuri.

\end{enumerate}




